\chapter{\IfLanguageName{dutch}{Methodologie}{Methodology}}%
\label{ch:methodologie}

In dit hoofdstuk wordt de onderzoeksopzet gedetailleerd toegelicht. Om te bepalen welke AI-geassisteerde workflow de meest kwalitatieve projectopstart oplevert, wordt een vergelijkend experiment uitgevoerd. Dit experiment toetst drie verschillende technologische benaderingen tegenover diverse inputbronnen en startcondities.

\section{Instrumentatie: Beschrijving van de Tools}
\label{sec:tools-beschrijving}

Voor dit onderzoek zijn drie categorieën van tools geselecteerd, elk met een verschillende graad van specialisatie en context-bewustzijn.

\begin{description}
    \item[Generieke LLM (ChatGPT/Claude):] Deze tools vertegenwoordigen de 'baseline'. Ze beschikken over een enorme hoeveelheid algemene programmeerkennis, maar hebben geen toegang tot de specifieke, niet-publieke componentenbibliotheek van de agency.
    \item[Gespecialiseerde UI-Generator (v0.dev):] v0 is een generatieve AI-tool van Vercel, specifiek geoptimaliseerd voor het bouwen van gebruikersinterfaces in React en Tailwind CSS. Deze tool is geselecteerd omdat het de huidige industriestandaard is voor snelle UI-prototyping.
    \item[Context-Aware setup (MCP-gebaseerde PoC):] Dit is de experimentele opstelling waarbij een LLM (zoals Claude 3.5 Sonnet) via het \textit{Model Context Protocol} (MCP) wordt verbonden met een lokale server. Deze server ontsluit de documentatie, code-standaarden en de React-componenten van Ballistix Digital.
\end{description}

\section{Experimentele Procedure en Workflow}
\label{sec:procedure}

Het experiment wordt uitgevoerd in een gecontroleerde omgeving om de variabelen zuiver te houden. De workflow verloopt voor elk scenario volgens een vast stramien:

\begin{enumerate}
    \item \textbf{Input-voorbereiding:} Voor elke testcase (visueel, functioneel, technisch) wordt een identiek pakket aan informatie klaargezet.
    \item \textbf{Promp-engineering:} Er wordt gebruikgemaakt van een gestandaardiseerd 'System Prompt' voor de generieke en MCP-tools om de instructie-bias te minimaliseren.
    \item \textbf{Generatie-fase:} De code wordt gegenereerd en direct geëxporteerd naar een specifieke Git-branch per scenario.
    \item \textbf{Integratie-test:} De gegenereerde code wordt geplaatst in zowel een leeg project als in de bestaande Ballistix-projecttemplate om de compatibiliteit te testen.
\end{enumerate}

\section{Onderzoeksmatrix}
\label{sec:testmatrix}

De kern van de methodologie is de test-matrix waarin de drie tools worden afgezet tegen drie kwalitatief verschillende inputbronnen. Dit resulteert in negen unieke testscenario's (zie Tabel~\ref{tab:testmatrix-uitgebreid}).

\begin{table}[h]
  \centering
  \begin{tabular}{l|ccc}
    \toprule
    \textbf{Inputbron} & \textbf{Generiek} & \textbf{v0.dev} & \textbf{MCP-PoC} \\
    \midrule
    \textbf{Visueel} (Screenshot/Figma) & Scenario A1 & Scenario A2 & Scenario A3 \\
    \textbf{Functioneel} (User Story) & Scenario B1 & Scenario B2 & Scenario B3 \\
    \textbf{Technisch} (Dev Prompt) & Scenario C1 & Scenario C2 & Scenario C3 \\
    \bottomrule
  \end{tabular}
  \caption{Uitgebreide test-matrix voor de vergelijking van AI-geassisteerde workflows.}
  \label{tab:testmatrix-uitgebreid}
\end{table}



\section{Kwaliteitsmeting en Evaluatiekader}
\label{sec:kwaliteitsmeting}

Om de "correctheid" van de oplossing objectief te bepalen, wordt elk scenario geëvalueerd op basis van vier kwantitatieve en kwalitatieve KPI's:

\subsection{1. Design System Compliance}
Er wordt gemeten in welke mate de AI gebruikmaakt van de bestaande \textit{Design Tokens} (kleuren, spatiëring) en herbruikbare React-componenten. 
\begin{itemize}
    \item \textbf{Score:} Percentage correct aangeroepen agency-componenten t.o.v. nieuw gegenereerde (generieke) HTML-elementen.
\end{itemize}

\subsection{2. Technical Correctness \& Compilation}
Dit criterium meet of de code direct uitvoerbaar is binnen de agency-stack.
\begin{itemize}
    \item \textbf{Meting:} Aantal \textit{Linting}-fouten en TypeScript-errors bij de eerste build.
\end{itemize}

\subsection{3. Correction Overhead (Tijdmeting)}
Dit is de meest kritieke factor voor de agency-efficiëntie. Er wordt gemeten hoeveel tijd een senior developer nodig heeft om de gegenereerde code "productierijp" te maken (refactoring).
\begin{itemize}
    \item \textbf{Meting:} Tijd in minuten benodigd voor correcties.
\end{itemize}

\subsection{4. Architectural Fit}
Er wordt beoordeeld of de AI de mappenstructuur en de scheiding van logica (components, hooks, types) respecteert zoals gedefinieerd in de Ballistix-template.
\begin{itemize}
    \item \textbf{Meting:} Score op een Likert-schaal (1-5) door een technisch expert.
\end{itemize}